% ---------------------------------------------------------------------------
% preamble.tex --- shared Beamer style for the four workshop decks
%
% Loaded by each deck with % ---------------------------------------------------------------------------
% preamble.tex --- shared Beamer style for the four workshop decks
%
% Loaded by each deck with % ---------------------------------------------------------------------------
% preamble.tex --- shared Beamer style for the four workshop decks
%
% Loaded by each deck with % ---------------------------------------------------------------------------
% preamble.tex --- shared Beamer style for the four workshop decks
%
% Loaded by each deck with \input{../common/preamble}. Same fonts, colours
% and bibliography settings as the book (book/preamble.tex): Latin Modern,
% black on white, one grey for code, one accent colour for what is untrusted
% or wrong. Built with pdflatex + biber, like the book.
% ---------------------------------------------------------------------------

\usepackage[T1]{fontenc}
\usepackage[utf8]{inputenc}
\usepackage{lmodern}
\usepackage{microtype}
\usepackage{booktabs}
\usepackage{tabularx}
\usepackage{listings}
\usepackage{xcolor}
\usepackage{csquotes}
\usepackage{graphicx}

% Figures: the book's figures as the 300 dpi PNGs the book itself includes
% (book/figures/png, tracked in git and current; figures/pdf is not rebuilt
% after every change), and the first pages of the papers
% (common/screenshots). Both are found by name without an extension.
\graphicspath{{../../book/figures/png/}{../common/screenshots/}}
\DeclareGraphicsExtensions{.png,.pdf}

% --- look ------------------------------------------------------------------
\usetheme{default}
\usefonttheme{serif}
\definecolor{accent}{HTML}{9B2D20}
\definecolor{codebg}{gray}{0.96}
\definecolor{rule}{gray}{0.75}
\definecolor{dim}{gray}{0.45}
\setbeamercolor{structure}{fg=black}
\setbeamercolor{normal text}{fg=black,bg=white}
\setbeamercolor{alerted text}{fg=accent}
\setbeamercolor{section in toc}{fg=black}
\setbeamercolor{footline}{fg=dim}
\setbeamerfont{title}{series=\bfseries,size=\Large}
\setbeamerfont{subtitle}{series=\mdseries,size=\normalsize}
\setbeamerfont{frametitle}{series=\bfseries,size=\large}
\setbeamerfont{footline}{size=\scriptsize}
\setbeamertemplate{navigation symbols}{}
\setbeamertemplate{itemize items}[circle]
\setbeamertemplate{enumerate items}[default]
\setbeamertemplate{caption}{\scriptsize\insertcaption}
\setbeamertemplate{frametitle}{%
  \vskip0.6em\usebeamerfont{frametitle}\insertframetitle\par
  \vskip0.2em{\color{rule}\hrule height 0.4pt}\vskip0.4em}
\setbeamertemplate{footline}{%
  \hbox to \paperwidth{\hspace{1.2em}\usebeamerfont{footline}\usebeamercolor[fg]{footline}%
    \insertshorttitle\ --- \insertsection\hfill
    \insertframenumber/\inserttotalframenumber\hspace{1.2em}}\vskip0.6em}
\setbeamersize{text margin left=1.4em,text margin right=1.4em}
% An outline frame at the start of each section, the current one in black.
\AtBeginSection[]{\begin{frame}{Outline}\tableofcontents[currentsection]\end{frame}}
% A frame that breaks over pages is numbered only from its second page.
\setbeamertemplate{frametitle continuation}[from second]

% --- code listings (as in the book) ----------------------------------------
\lstset{
  basicstyle=\ttfamily\footnotesize,
  backgroundcolor=\color{codebg},
  frame=none,
  breaklines=true,
  breakatwhitespace=true,
  columns=fullflexible,
  keepspaces=true,
  upquote=true,
  showstringspaces=false,
  xleftmargin=0.6em,
  framexleftmargin=0.6em,
  aboveskip=0.5em,
  belowskip=0.5em,
  literate={~}{{\textasciitilde}}1
           {>}{{\textgreater}}1
           {<}{{\textless}}1
           {ö}{{\"o}}1 {ü}{{\"u}}1 {é}{{\'e}}1,
}
\lstdefinestyle{shell}{language=bash}

% --- bibliography (same settings as the book) ------------------------------
\usepackage[
  backend=biber,
  style=authoryear,
  sorting=nyt,
  maxbibnames=6,
  minbibnames=3,
  maxcitenames=2,
  giveninits=true,
  uniquename=init,
  doi=true,
  url=true,
  isbn=false,
  eprint=true,
]{biblatex}
\addbibresource{../../book/references.bib}
\renewcommand*{\bibfont}{\scriptsize}
\setlength{\bibitemsep}{0.25em}
\setbeamertemplate{bibliography item}{}   % no icon before each entry

% --- shorthand -------------------------------------------------------------
\newcommand{\code}[1]{\texttt{#1}}
\newcommand{\term}[1]{\emph{#1}}
% A key point, as in the book: no box.
\newcommand{\kp}[1]{\par\medskip{\bfseries Key point.}\enspace\emph{#1}\par}
% \fig[height]{name}: one of the book's figures, scaled to fit the frame.
\newcommand{\fig}[2][0.72\textheight]{%
  \centering\includegraphics[height=#1,width=\textwidth,keepaspectratio]{#2}\par}
% \screenshot{key}: the first page of a paper, with a thin frame.
\newcommand{\screenshot}[2][0.66\textheight]{%
  \setlength{\fboxsep}{0pt}\setlength{\fboxrule}{0.3pt}%
  \fbox{\includegraphics[height=#1,width=\linewidth,keepaspectratio]{#2}}}
% \paperslide{key}{text}: screenshot on the left, what the book takes from the
% paper on the right. Use inside a frame; add \source after it.
\newcommand{\paperslide}[2]{%
  \begin{columns}[T]
    \begin{column}{0.40\textwidth}\centering\screenshot{#1}\end{column}
    \begin{column}{0.58\textwidth}\small #2\end{column}
  \end{columns}}
% Notes in references.bib record how a web source was saved and dated; they
% belong in the reference list, not on a \source line, so cite commands
% (including \fullcite) leave them out.
\AtEveryCitekey{\clearfield{note}}
% \source{...}: a small grey line at the foot of a frame, for the source.
\newcommand{\source}[1]{\vskip0pt plus 1filll\par{\scriptsize\color{dim}#1\par}}
. Same fonts, colours
% and bibliography settings as the book (book/preamble.tex): Latin Modern,
% black on white, one grey for code, one accent colour for what is untrusted
% or wrong. Built with pdflatex + biber, like the book.
% ---------------------------------------------------------------------------

\usepackage[T1]{fontenc}
\usepackage[utf8]{inputenc}
\usepackage{lmodern}
\usepackage{microtype}
\usepackage{booktabs}
\usepackage{tabularx}
\usepackage{listings}
\usepackage{xcolor}
\usepackage{csquotes}
\usepackage{graphicx}

% Figures: the book's figures as the 300 dpi PNGs the book itself includes
% (book/figures/png, tracked in git and current; figures/pdf is not rebuilt
% after every change), and the first pages of the papers
% (common/screenshots). Both are found by name without an extension.
\graphicspath{{../../book/figures/png/}{../common/screenshots/}}
\DeclareGraphicsExtensions{.png,.pdf}

% --- look ------------------------------------------------------------------
\usetheme{default}
\usefonttheme{serif}
\definecolor{accent}{HTML}{9B2D20}
\definecolor{codebg}{gray}{0.96}
\definecolor{rule}{gray}{0.75}
\definecolor{dim}{gray}{0.45}
\setbeamercolor{structure}{fg=black}
\setbeamercolor{normal text}{fg=black,bg=white}
\setbeamercolor{alerted text}{fg=accent}
\setbeamercolor{section in toc}{fg=black}
\setbeamercolor{footline}{fg=dim}
\setbeamerfont{title}{series=\bfseries,size=\Large}
\setbeamerfont{subtitle}{series=\mdseries,size=\normalsize}
\setbeamerfont{frametitle}{series=\bfseries,size=\large}
\setbeamerfont{footline}{size=\scriptsize}
\setbeamertemplate{navigation symbols}{}
\setbeamertemplate{itemize items}[circle]
\setbeamertemplate{enumerate items}[default]
\setbeamertemplate{caption}{\scriptsize\insertcaption}
\setbeamertemplate{frametitle}{%
  \vskip0.6em\usebeamerfont{frametitle}\insertframetitle\par
  \vskip0.2em{\color{rule}\hrule height 0.4pt}\vskip0.4em}
\setbeamertemplate{footline}{%
  \hbox to \paperwidth{\hspace{1.2em}\usebeamerfont{footline}\usebeamercolor[fg]{footline}%
    \insertshorttitle\ --- \insertsection\hfill
    \insertframenumber/\inserttotalframenumber\hspace{1.2em}}\vskip0.6em}
\setbeamersize{text margin left=1.4em,text margin right=1.4em}
% An outline frame at the start of each section, the current one in black.
\AtBeginSection[]{\begin{frame}{Outline}\tableofcontents[currentsection]\end{frame}}
% A frame that breaks over pages is numbered only from its second page.
\setbeamertemplate{frametitle continuation}[from second]

% --- code listings (as in the book) ----------------------------------------
\lstset{
  basicstyle=\ttfamily\footnotesize,
  backgroundcolor=\color{codebg},
  frame=none,
  breaklines=true,
  breakatwhitespace=true,
  columns=fullflexible,
  keepspaces=true,
  upquote=true,
  showstringspaces=false,
  xleftmargin=0.6em,
  framexleftmargin=0.6em,
  aboveskip=0.5em,
  belowskip=0.5em,
  literate={~}{{\textasciitilde}}1
           {>}{{\textgreater}}1
           {<}{{\textless}}1
           {ö}{{\"o}}1 {ü}{{\"u}}1 {é}{{\'e}}1,
}
\lstdefinestyle{shell}{language=bash}

% --- bibliography (same settings as the book) ------------------------------
\usepackage[
  backend=biber,
  style=authoryear,
  sorting=nyt,
  maxbibnames=6,
  minbibnames=3,
  maxcitenames=2,
  giveninits=true,
  uniquename=init,
  doi=true,
  url=true,
  isbn=false,
  eprint=true,
]{biblatex}
\addbibresource{../../book/references.bib}
\renewcommand*{\bibfont}{\scriptsize}
\setlength{\bibitemsep}{0.25em}
\setbeamertemplate{bibliography item}{}   % no icon before each entry

% --- shorthand -------------------------------------------------------------
\newcommand{\code}[1]{\texttt{#1}}
\newcommand{\term}[1]{\emph{#1}}
% A key point, as in the book: no box.
\newcommand{\kp}[1]{\par\medskip{\bfseries Key point.}\enspace\emph{#1}\par}
% \fig[height]{name}: one of the book's figures, scaled to fit the frame.
\newcommand{\fig}[2][0.72\textheight]{%
  \centering\includegraphics[height=#1,width=\textwidth,keepaspectratio]{#2}\par}
% \screenshot{key}: the first page of a paper, with a thin frame.
\newcommand{\screenshot}[2][0.66\textheight]{%
  \setlength{\fboxsep}{0pt}\setlength{\fboxrule}{0.3pt}%
  \fbox{\includegraphics[height=#1,width=\linewidth,keepaspectratio]{#2}}}
% \paperslide{key}{text}: screenshot on the left, what the book takes from the
% paper on the right. Use inside a frame; add \source after it.
\newcommand{\paperslide}[2]{%
  \begin{columns}[T]
    \begin{column}{0.40\textwidth}\centering\screenshot{#1}\end{column}
    \begin{column}{0.58\textwidth}\small #2\end{column}
  \end{columns}}
% Notes in references.bib record how a web source was saved and dated; they
% belong in the reference list, not on a \source line, so cite commands
% (including \fullcite) leave them out.
\AtEveryCitekey{\clearfield{note}}
% \source{...}: a small grey line at the foot of a frame, for the source.
\newcommand{\source}[1]{\vskip0pt plus 1filll\par{\scriptsize\color{dim}#1\par}}
. Same fonts, colours
% and bibliography settings as the book (book/preamble.tex): Latin Modern,
% black on white, one grey for code, one accent colour for what is untrusted
% or wrong. Built with pdflatex + biber, like the book.
% ---------------------------------------------------------------------------

\usepackage[T1]{fontenc}
\usepackage[utf8]{inputenc}
\usepackage{lmodern}
\usepackage{microtype}
\usepackage{booktabs}
\usepackage{tabularx}
\usepackage{listings}
\usepackage{xcolor}
\usepackage{csquotes}
\usepackage{graphicx}

% Figures: the book's figures as the 300 dpi PNGs the book itself includes
% (book/figures/png, tracked in git and current; figures/pdf is not rebuilt
% after every change), and the first pages of the papers
% (common/screenshots). Both are found by name without an extension.
\graphicspath{{../../book/figures/png/}{../common/screenshots/}}
\DeclareGraphicsExtensions{.png,.pdf}

% --- look ------------------------------------------------------------------
\usetheme{default}
\usefonttheme{serif}
\definecolor{accent}{HTML}{9B2D20}
\definecolor{codebg}{gray}{0.96}
\definecolor{rule}{gray}{0.75}
\definecolor{dim}{gray}{0.45}
\setbeamercolor{structure}{fg=black}
\setbeamercolor{normal text}{fg=black,bg=white}
\setbeamercolor{alerted text}{fg=accent}
\setbeamercolor{section in toc}{fg=black}
\setbeamercolor{footline}{fg=dim}
\setbeamerfont{title}{series=\bfseries,size=\Large}
\setbeamerfont{subtitle}{series=\mdseries,size=\normalsize}
\setbeamerfont{frametitle}{series=\bfseries,size=\large}
\setbeamerfont{footline}{size=\scriptsize}
\setbeamertemplate{navigation symbols}{}
\setbeamertemplate{itemize items}[circle]
\setbeamertemplate{enumerate items}[default]
\setbeamertemplate{caption}{\scriptsize\insertcaption}
\setbeamertemplate{frametitle}{%
  \vskip0.6em\usebeamerfont{frametitle}\insertframetitle\par
  \vskip0.2em{\color{rule}\hrule height 0.4pt}\vskip0.4em}
\setbeamertemplate{footline}{%
  \hbox to \paperwidth{\hspace{1.2em}\usebeamerfont{footline}\usebeamercolor[fg]{footline}%
    \insertshorttitle\ --- \insertsection\hfill
    \insertframenumber/\inserttotalframenumber\hspace{1.2em}}\vskip0.6em}
\setbeamersize{text margin left=1.4em,text margin right=1.4em}
% An outline frame at the start of each section, the current one in black.
\AtBeginSection[]{\begin{frame}{Outline}\tableofcontents[currentsection]\end{frame}}
% A frame that breaks over pages is numbered only from its second page.
\setbeamertemplate{frametitle continuation}[from second]

% --- code listings (as in the book) ----------------------------------------
\lstset{
  basicstyle=\ttfamily\footnotesize,
  backgroundcolor=\color{codebg},
  frame=none,
  breaklines=true,
  breakatwhitespace=true,
  columns=fullflexible,
  keepspaces=true,
  upquote=true,
  showstringspaces=false,
  xleftmargin=0.6em,
  framexleftmargin=0.6em,
  aboveskip=0.5em,
  belowskip=0.5em,
  literate={~}{{\textasciitilde}}1
           {>}{{\textgreater}}1
           {<}{{\textless}}1
           {ö}{{\"o}}1 {ü}{{\"u}}1 {é}{{\'e}}1,
}
\lstdefinestyle{shell}{language=bash}

% --- bibliography (same settings as the book) ------------------------------
\usepackage[
  backend=biber,
  style=authoryear,
  sorting=nyt,
  maxbibnames=6,
  minbibnames=3,
  maxcitenames=2,
  giveninits=true,
  uniquename=init,
  doi=true,
  url=true,
  isbn=false,
  eprint=true,
]{biblatex}
\addbibresource{../../book/references.bib}
\renewcommand*{\bibfont}{\scriptsize}
\setlength{\bibitemsep}{0.25em}
\setbeamertemplate{bibliography item}{}   % no icon before each entry

% --- shorthand -------------------------------------------------------------
\newcommand{\code}[1]{\texttt{#1}}
\newcommand{\term}[1]{\emph{#1}}
% A key point, as in the book: no box.
\newcommand{\kp}[1]{\par\medskip{\bfseries Key point.}\enspace\emph{#1}\par}
% \fig[height]{name}: one of the book's figures, scaled to fit the frame.
\newcommand{\fig}[2][0.72\textheight]{%
  \centering\includegraphics[height=#1,width=\textwidth,keepaspectratio]{#2}\par}
% \screenshot{key}: the first page of a paper, with a thin frame.
\newcommand{\screenshot}[2][0.66\textheight]{%
  \setlength{\fboxsep}{0pt}\setlength{\fboxrule}{0.3pt}%
  \fbox{\includegraphics[height=#1,width=\linewidth,keepaspectratio]{#2}}}
% \paperslide{key}{text}: screenshot on the left, what the book takes from the
% paper on the right. Use inside a frame; add \source after it.
\newcommand{\paperslide}[2]{%
  \begin{columns}[T]
    \begin{column}{0.40\textwidth}\centering\screenshot{#1}\end{column}
    \begin{column}{0.58\textwidth}\small #2\end{column}
  \end{columns}}
% Notes in references.bib record how a web source was saved and dated; they
% belong in the reference list, not on a \source line, so cite commands
% (including \fullcite) leave them out.
\AtEveryCitekey{\clearfield{note}}
% \source{...}: a small grey line at the foot of a frame, for the source.
\newcommand{\source}[1]{\vskip0pt plus 1filll\par{\scriptsize\color{dim}#1\par}}
. Same fonts, colours
% and bibliography settings as the book (book/preamble.tex): Latin Modern,
% black on white, one grey for code, one accent colour for what is untrusted
% or wrong. Built with pdflatex + biber, like the book.
% ---------------------------------------------------------------------------

\usepackage[T1]{fontenc}
\usepackage[utf8]{inputenc}
\usepackage{lmodern}
\usepackage{microtype}
\usepackage{booktabs}
\usepackage{tabularx}
\usepackage{listings}
\usepackage{xcolor}
\usepackage{csquotes}
\usepackage{graphicx}

% Figures: the book's figures as the 300 dpi PNGs the book itself includes
% (book/figures/png, tracked in git and current; figures/pdf is not rebuilt
% after every change), and the first pages of the papers
% (common/screenshots). Both are found by name without an extension.
\graphicspath{{../../book/figures/png/}{../common/screenshots/}}
\DeclareGraphicsExtensions{.png,.pdf}

% --- look ------------------------------------------------------------------
\usetheme{default}
\usefonttheme{serif}
\definecolor{accent}{HTML}{9B2D20}
\definecolor{codebg}{gray}{0.96}
\definecolor{rule}{gray}{0.75}
\definecolor{dim}{gray}{0.45}
\setbeamercolor{structure}{fg=black}
\setbeamercolor{normal text}{fg=black,bg=white}
\setbeamercolor{alerted text}{fg=accent}
\setbeamercolor{section in toc}{fg=black}
\setbeamercolor{footline}{fg=dim}
\setbeamerfont{title}{series=\bfseries,size=\Large}
\setbeamerfont{subtitle}{series=\mdseries,size=\normalsize}
\setbeamerfont{frametitle}{series=\bfseries,size=\large}
\setbeamerfont{footline}{size=\scriptsize}
\setbeamertemplate{navigation symbols}{}
\setbeamertemplate{itemize items}[circle]
\setbeamertemplate{enumerate items}[default]
\setbeamertemplate{caption}{\scriptsize\insertcaption}
\setbeamertemplate{frametitle}{%
  \vskip0.6em\usebeamerfont{frametitle}\insertframetitle\par
  \vskip0.2em{\color{rule}\hrule height 0.4pt}\vskip0.4em}
\setbeamertemplate{footline}{%
  \hbox to \paperwidth{\hspace{1.2em}\usebeamerfont{footline}\usebeamercolor[fg]{footline}%
    \insertshorttitle\ --- \insertsection\hfill
    \insertframenumber/\inserttotalframenumber\hspace{1.2em}}\vskip0.6em}
\setbeamersize{text margin left=1.4em,text margin right=1.4em}
% An outline frame at the start of each section, the current one in black.
\AtBeginSection[]{\begin{frame}{Outline}\tableofcontents[currentsection]\end{frame}}
% A frame that breaks over pages is numbered only from its second page.
\setbeamertemplate{frametitle continuation}[from second]

% --- code listings (as in the book) ----------------------------------------
\lstset{
  basicstyle=\ttfamily\footnotesize,
  backgroundcolor=\color{codebg},
  frame=none,
  breaklines=true,
  breakatwhitespace=true,
  columns=fullflexible,
  keepspaces=true,
  upquote=true,
  showstringspaces=false,
  xleftmargin=0.6em,
  framexleftmargin=0.6em,
  aboveskip=0.5em,
  belowskip=0.5em,
  literate={~}{{\textasciitilde}}1
           {>}{{\textgreater}}1
           {<}{{\textless}}1
           {ö}{{\"o}}1 {ü}{{\"u}}1 {é}{{\'e}}1,
}
\lstdefinestyle{shell}{language=bash}

% --- bibliography (same settings as the book) ------------------------------
\usepackage[
  backend=biber,
  style=authoryear,
  sorting=nyt,
  maxbibnames=6,
  minbibnames=3,
  maxcitenames=2,
  giveninits=true,
  uniquename=init,
  doi=true,
  url=true,
  isbn=false,
  eprint=true,
]{biblatex}
\addbibresource{../../book/references.bib}
\renewcommand*{\bibfont}{\scriptsize}
\setlength{\bibitemsep}{0.25em}
\setbeamertemplate{bibliography item}{}   % no icon before each entry

% --- shorthand -------------------------------------------------------------
\newcommand{\code}[1]{\texttt{#1}}
\newcommand{\term}[1]{\emph{#1}}
% A key point, as in the book: no box.
\newcommand{\kp}[1]{\par\medskip{\bfseries Key point.}\enspace\emph{#1}\par}
% \fig[height]{name}: one of the book's figures, scaled to fit the frame.
\newcommand{\fig}[2][0.72\textheight]{%
  \centering\includegraphics[height=#1,width=\textwidth,keepaspectratio]{#2}\par}
% \screenshot{key}: the first page of a paper, with a thin frame.
\newcommand{\screenshot}[2][0.66\textheight]{%
  \setlength{\fboxsep}{0pt}\setlength{\fboxrule}{0.3pt}%
  \fbox{\includegraphics[height=#1,width=\linewidth,keepaspectratio]{#2}}}
% \paperslide{key}{text}: screenshot on the left, what the book takes from the
% paper on the right. Use inside a frame; add \source after it.
\newcommand{\paperslide}[2]{%
  \begin{columns}[T]
    \begin{column}{0.40\textwidth}\centering\screenshot{#1}\end{column}
    \begin{column}{0.58\textwidth}\small #2\end{column}
  \end{columns}}
% Notes in references.bib record how a web source was saved and dated; they
% belong in the reference list, not on a \source line, so cite commands
% (including \fullcite) leave them out.
\AtEveryCitekey{\clearfield{note}}
% \source{...}: a small grey line at the foot of a frame, for the source.
\newcommand{\source}[1]{\vskip0pt plus 1filll\par{\scriptsize\color{dim}#1\par}}
